\subsection{The second real embedding in the estimate at primes with selected orbits}
\label{subsec:second-period}

The first embedding supplies a gain from the selected orbits. The product estimate also needs the contribution at the other real embedding.

Fix a word, a used tuple, and a coefficient place over a prime
\(p\) of the selected class. The tuple's twist support avoids \(p\).
The CM field splits at the places of its real subfield used for the
local matrix pairs, while \(p\ge5\) is inert in the quadratic
field of the elliptic curve. The real embeddings are indexed by \(\{0,1\}\times H\), where \(H\) is the finite abelian factor indexing the second coordinate. Writing \(h=h_p\), its Frobenius action on the labels satisfies
\[
       (i,a)\longmapsto(i,a+h),\qquad d=\operatorname{ord}(h)>0,\quad 2\mid d .
\] Let \(R\)
contain one representative of each coset modulo \(\langle h\rangle\).
The choice of first-embedding places in \eqref{eq:tag-p7} assigns to each
\(r\in R\) the cycle starting at \((0,r)\).

\begin{proposition}[The complementary contribution, including class denominators]\label{prop:complementary-period-price}
In the selected class and on the indicated level-free tuple,
the complementary lower bounds \(L_j\) at the successive stages satisfy
\[
 \sum_{j=1}^m L_j\ge
 {\bf1}_{p\le X''(n)}\sum_{j=1}^m
   \bigl(F_{\rm inert}(j,p)-E_{\rm inert}(j,p)\bigr)
  -(m-1)v_p(D_{\Gamma,u}).
\]
The terminal contraction may cost at most one more
\(v_p(D_{\Gamma,u})\).
\end{proposition}

\begin{proof}
\proofstep{Assign each complementary cycle in the fixed embedding coordinates.}
The comparison of embeddings identifies the pairs \((v,s)\) of
a local place and one of its residue embeddings bijectively with
the real positions. In fact the position of the first embedding
of \(v\) determines \((b_v,t_v)\). Composing with residue
Frobenius gives \((b_v,t_v+s h)\) at embedding \(s\).
The residue degree is the order of the same Frobenius element in
the unramified host. Thus the map gives each position of this
cycle exactly once. A place outside the chosen first-embedding
places must have \(b_v=1\), since each \((0,r)\) already occurs.

For each \(r\in R\), the unique preimage of \((1,r)\) gives its
second-embedding place. The latter is independent of the starting
letter in the coset. In particular, writing \(V\) for the chosen
first-embedding places, we have
\[
  \sum_{v\notin V} f(t_v)=\sum_{r\in R}f(r)
  \qquad\bigl(f(t+h)=f(t)\bigr).
  \tag{S1}\label{eq:tag-s1}
\]
The set \(V\) contains the first-embedding places even for
orbits at which the generator uses only the scalar option.

Consider a transition of ranks \(N=N_{j-1},N'=N_j\).
For a second-embedding cycle let \(w_s\) be its selected positive
stage weights. Use the linear program in
Section~\ref{subsec:cycle-jets} and put
\[
 F_{jr}={\cal F}_p(((k-1)w_s)_s),\qquad N_* = N_m .
\]
There are no designated tops at this real embedding. The bounds
on the shape of the stage diagram here read
\[
 H_s=(k-1)w_s+e_s,\qquad 0\le e_s\le4w_s,\qquad H_s\le a^2 .
\]
Here \(H_s\) counts all the boxes. In particular the list of
distinguished rows reserved in Proposition~\ref{prop:cyclic-jet}
is empty.

For a monomial in the integral parameters, the trace
basis substitution of Section~\ref{subsec:curve-exponents} and the
flow comparison \eqref{eq:flow-comparison} bound its normalized
differential coefficient from below by
\[
 F_{jr}-11\sum_s w_s/p-dNN'(1+\log_p a^2).
\]
To see how the two logarithms bound the error, on \(p\le X''(n)\le a\) set
\[
 E_{jr}= c_6\sum_s w_s/p
 +dNN'\bigl(1+\log_p(4dN_*a)\bigr)
 +d\,d_1 N_*^2\log_p(2a).
 \tag{S2}\label{eq:tag-s2}
\]
Use \(c_6=12,d_1=5\). The last two terms majorize
\(dNN'(1+2\log_p a)\), since \(N,N'\le N_*\).
In the general orbit formula both the maximum-width loss and
the squared digit loss have a factor counting tops. This count
is zero on the cycles in \eqref{eq:tag-s2}.

For precision about coefficients, take a monomial in the integral
formal parameters over the Witt ring. The differential polynomial
has its coefficients in the Witt fraction field after the trace
basis change for the local model. Its normalized valuation there
is integer-valued. For each coefficient \(b\) in this unramified valuation, integrality of its order strengthens the given lower bound to
\[
  v_p(b)\ge F_{jr}-E_{jr}\quad\Longrightarrow\quad b\in p^{\lceil F_{jr}-E_{jr}\rceil}W.
\]
For a negative exponent the right side denotes the corresponding fractional ideal. An integral multiplier in a larger coefficient
ring gives the same bound. Indeed truncate its expansion in
parameters at the total degree of the differential operator.
Substitution by series with zero constant term preserves the
order filtration. The remainder still has order strictly greater
than the differential degree after multiplication by the phase,
whose constant coefficient is one.

For a differential operator of total order \(h\), truncation is justified by
\[
  F\in (X)^{h+1},\quad |\alpha|\le h\quad\Longrightarrow\quad \partial^\alpha F(0)=0 .
\]
 The finitely many surviving
terms are the Witt monomial evaluations multiplied by their
integral coefficients. This proves the bound for arbitrary
integral multipliers, including those involving variables at
other places. Above the cutoff one only uses the bound zero,
which follows from the integral ordinary derivatives of the
phase and of the integral multiplier as in
Section~\ref{subsec:hecke-entries}.

Both \(F_{jr}\) and \(E_{jr}\) are invariant under rotating the
cycle. Rotation bijects feasible flows, preserving the objective,
and permutes the weights in \eqref{eq:tag-s2}. The full inert forced value
\(F_{\rm inert}(j,p)\) is the sum of \(F\) over both sets of
cycles. Equivalently it is the sum over positions divided by \(d\).
Write \(F_{j,0}\) for the sum over first-embedding cycles. In this
prime class the inert error in the division table is
\[
 E_{\rm inert}(j,p)=\max(0,F_{j,0})+\sum_{r\in R}E_{jr}.
\]
Indeed the first term of that table is the positive part of the
sum at first positions divided by \(d\). The other terms sum
\eqref{eq:tag-s2} over the second positions, since each cycle has \(d\)
positions. The exceptional term for a nonzero step of odd order
vanishes. It follows from \eqref{eq:tag-s1} that for \(p\le X''(n)\)
\[
 \sum_{v\notin V}(F_{j,t_v}-E_{j,t_v})
 \ \ge\ F_{\rm inert}(j,p)-E_{\rm inert}(j,p).
 \tag{S3}\label{eq:tag-s3}
\]
This inequality subtracts the positive part after summing over
the first cycles. It holds for either sign of each forced value.

Finally let \(d_i\) be the maximum valuation of the point
stabilizer orders at the source of transition \(i+1\), with an
empty maximum taken to be zero. These orders divide the common
integer \(D_{\Gamma,u}\). At level zero the two fine type levels
have trivial rational stabilizer by
Section~\ref{sec:ramified-column}. \proofstep{Charge only the incoming class denominators, then include the terminal contraction.}
The class sum at each next
transition costs at most \(d_i\), by the ultrametric inequality.
Consequently \(\sum_{i=0}^{m-1}d_i\le(m-1)v_p(D_{\Gamma,u})\).
Substituting \eqref{eq:tag-s3} into \eqref{eq:tag-p7} gives
\[
 \sum_{j=1}^m L_j\ \ge {\bf1}_{p\le X''(n)}\sum_{j=1}^m
      (F_{\rm inert}(j,p)-E_{\rm inert}(j,p))
       -(m-1)v_p(D_{\Gamma,u}).
 \tag{S4}\label{eq:tag-s4}
\]
The terminal period contraction in \eqref{eq:terminal-period-valuation}
can consume the remaining copy of \(v_p(D_{\Gamma,u})\). None
of these losses depends on the number of paths or on the
degree of the field of coefficients. \label{cor:complementary-period-1}
\end{proof}

